\section{Related Work}

\para{Weight-Space Ensembles.}
Traditional ensembling averages the outputs of multiple independently trained models in logit space. Recently, weight averaging and weight-space ensembles have gained traction due to their computational efficiency at inference time. \citet{wortsman2022model} introduced Model Soups, showing that averaging the weights of models fine-tuned with different hyperparameters can significantly improve performance on target distributions without increasing inference cost. Building on this, \citet{gan2026neural} proposed the Neural Thicket hypothesis and the \randopt method, demonstrating that dense neighborhoods of high-performing task experts exist around pretrained weights, and can be discovered via random Gaussian perturbations. Unlike these works, which primarily focus on clean benchmark accuracy, we explicitly investigate the robustness of these weight-space neighborhoods to adversarial and distribution shifts.

\para{Robustness of Interpolated Models.}
A growing body of literature explores how combining models affects their resilience to distribution shifts. \citet{croce2023seasoning} demonstrated that ``seasoning'' model soups by interpolating models fine-tuned with different adversarial objectives can smoothly trade off and improve robustness to multiple threats. However, their approach requires expensive adversarial fine-tuning of multiple models. In contrast, our work investigates whether robustness can be achieved purely through random and explicitly diverse perturbations around a single pretrained model, without the need for specialized adversarial training. We show that while standard random perturbations are insufficient and can even harm \ood robustness, explicitly enforcing diversity via \methodname recovers and enhances ensemble resilience.